
With a complete representation of the search space by relational algebra, our
next step is to explore this space and find the optimal expression in it.
Traditional optimizing compilers commonly resort to heuristics to select from
available rewrites to apply. SystemML implements a number of heuristics for its
algebraic rewrite rules, and we discuss a few categories of them here.

\textsc{Competing or Conflicting Rewrites} The same expression may be eligible
for more than one rewrites. For example, $sum(AB)$ rewrites to
$sum(sum_{col}(A)^T*sum_{row}(B))$, but when both $A$ and $B$ are vectors the
expression can also be rewritten to a single dot product. SystemML then
implements heuristics to only perform the first rewrite when the expression is
not a dot product. In the worst case, a set of rules interacting with each other
may create a quadratic number of such conflicts, complicating the codebase.

\textsc{Order of Rewrites} Some rewrite should be applied after others to be
effective. For example, $X/y$ could be rewritten to $X*1/y$ which may be more
efficient, since SystemML provides efficient implementation for sparse
multiplication but not for division. This rewrite should occur before constant
folding; otherwise it may create spurious expressions like $X / (1/y)
\rightarrow X * (1/(1/y))$, and without constant folding the double division
will persist. However, a rewrite like $1/(1 + exp(-X)) \rightarrow sigmoid(X)$
should come after constant folding, in order to cover expressions like $(3-2)/(1
+ exp(-X))$. Since SystemML requires all rewrites to happen in one phase and
constant folding another, it has to leave out\footnote{Another reason to leave
  out this rewrite is that $X*1/y$ rounds twice, whereas $X/y$ only rounds
  once.} rewrites like $X/y \rightarrow X*1/y$.

\textsc{Dependency on Input / Program Properties} Our example optimization from 
$sum((X-UV^T)^2)$  to $sum(X^2) -2U^TXV + U^TU*V^TV$
improves performance only if $X$ is sparse. Otherwise,
computing $X^2$ and $X*UV^T$ would both create dense intermediates. Similarly,
some rewrites depend on program properties like common subexpressions. Usually,
these rewrites only apply when the matched expression shares no CSE with others
in order to leverage common subexpression elimination. Testing input and program
properties like this becomes boilerplate code, making implementation tedious and
adds burden to maintenance.

\textsc{Composing Rewrites} Even more relevant to us is the problem of composing
larger rewrites out of smaller ones. Our equality rules $R_{EQ}$ are very
fine-grained, and any rule is unlikely to improve performance on its own. 
Our example optimization from $sum((X-UV^T)^2)$ to $sum(X^2) - 2U^TXV + U^TU * V^TV$ 
takes around 10 applications of $R_{EQ}$ rules. 
 If an optimizer applies rewrites one by one, it is
then very difficult, if not impossible, for it to discover the correct sequence
of rewrites that compose together and lead to the best performance.

Stepping back, the challenge of orchestrating rewrites is known as the
\emph{phase-ordering problem} in compiler optimization. Tate
et al. \cite{DBLP:journals/corr/abs-1012-1802} proposed a solution to this
problem dubbed
\textit{equality saturation}, which we adapt and extend in SPORES.

\subsection{Equality Saturation}

Equality saturation optimizes an expression in two steps:

\textit{Saturation}: given the input expression, the optimizer  enumerates
equivalent expressions and collects them into a compact representation called
the E-Graph \cite{10.5555/909447}.

\textit{Extraction}: given a cost function, the optimizer selects the optimal
expression from the E-Graph. An expression is represented by a subgraph of the
E-Graph, and the optimizer uses a constraint solver to find the subgraph that is
equivalent to the input, and is optimal according to the cost function.

\subsubsection*{The E-Graph Data Structure}

\begin{figure}
\centering
    \includegraphics[scale=0.6]{XYY} \qquad \includegraphics[scale=0.6]{assoc}
    \caption{E-Graph representing $(X*Y)*Y$ (left), and the graph after applying
      associativity to the root (right). New nodes are in gray. Each dashed box is an E-Class.}
    \label{assoc}
\end{figure}

An E-Graph represents sets of equivalent expressions. A node in the graph is
called an E-Class, which contains the root operators of a set of equivalent
expressions. The edges are similar to the edges in an abstract syntax tree; but
instead of pointing from an operator directly to a child, each edge points from
an operator to an E-Class of expressions. For example, in Figure~\ref{assoc} the
top class in the middle represents the set of equivalent expressions \{$(X*Y)*Y,
X*(Y*Y)$\}. Note that the class represents two expressions, each with 2
appearances of $Y$ and one appearance of $X$, whereas each variable only appears
once in the E-Graph. This is because the E-Graph makes sure its expressions share
all possible common subexpressions. As the size of the graph grows, this
compression becomes more and more notable; in some cases a graph can 
represent a number of expressions exponential to its
size \cite{DBLP:journals/corr/abs-1012-1802}. We take advantage of this
compression in SPORES to efficiently cover vast portions of the search
space. If saturation, as described below, carries out to convergence, the E-Graph represents the search space exhaustively. 

An E-Graph can also be seen as an AND-OR DAG over expressions. Each E-Class is an OR node whose children are equivalent expressions from
which the optimizer chooses from. Each operator is an AND node whose children must all be picked if the operator itself is picked. In this paper we favor the terms E-Graph and E-Class to emphasize each OR node is an equivalence class. 

\subsubsection*{Saturating the E-Graph}
At the beginning of the optimization process, the optimizer instantiates the
graph by inserting the nodes in the syntax tree of the input expression one by
one in post order. For example, for input $(X*Y)*Y$, we construct the left graph
in Figure~\ref{assoc} bottom-up. By inserting in post order, we readily exploit
existing common subexpressions in the input. Once the entire input expression is
inserted, the optimizer starts to extend the graph with new expressions equivalent to
the input. It considers a list of equations, and matches either side of the
equation to subgraphs of the E-Graph. If an equation matches, the optimizer then
inserts the expression on the other side of the equation to the graph. For
example, applying the associative rule extends the left graph in
Figure~\ref{assoc} with $X*(Y*Y)$, resulting in the right graph.
Figure~\ref{eqsat} shows the pseudo code for this process. While inserting new
expressions, the optimizer checks if any subexpression of the new expression is
already in the graph. If so, it reuses the existing node, thereby exploiting all
possible common-subexpressions to keep the E-Graph compact. In
Figure~\ref{assoc}, only two $*$ are added since the variables $X$ and $Y$ are
already in the graph. Once the entire new expression has been added, the
optimizer then merges the newly created E-Class at its root with the E-Class
containing the matched expression, asserting them equal. Importantly, the
optimizer also propagates the congruent closure of this new equality. For
example, when $A+A$ is merged with $2*A$, the optimizer also merges $(A+A)^2$
with $(2*A)^2$. Figure~\ref{fig:egraphadd} shows the pseudo code for adding an
expression to E-Graph. This process of match-and-insert is repeated until the
graph stops changing, or reaching a user-specified bound on the number of 
saturation iterations. If this process does converge, that means no rule can add new
expressions to the graph any more. If the set of rules are complete, as is our
$R_{EQ}$, convergence of saturation implies the resulting E-Graph represents
the transitive closure of the equality rules applied to the initial expression.
In other words, it contains \textit{all} expressions equivalent to the input under the equality rules.

\begin{figure}
\centering
\begin{minipage}{0.5\textwidth}
\begin{lstlisting}[language=Python]
def saturate(egraph, equations):
  for eq in equations: 
    matches = egraph.match(eq.lhs)
    for eclass in matches: 
      ec = egraph.add(eq.rhs)
      egraph.merge(eclass, c)
\end{lstlisting}
\end{minipage}
    \caption{Pseudo code for saturating the E-Graph given a set of equalities.
      \texttt{match} returns the IDs of the root class of any matching subgraph;
      \texttt{merge} combines two E-Classes given their IDs; it also propagates
      the congruent closure of the new equality.
      Figure~\ref{fig:egraphadd} defines \texttt{add}.}
    \label{eqsat}
\end{figure}
\begin{figure}
\centering
\begin{minipage}{0.5\textwidth}
\begin{lstlisting}[language=Python]
def add(expr):
  ID = egraph.find(expr)
  if ID != NULL:
    return ID
  else:
    cids = expr.children.map(add)
    ID = egraph.insert(expr.op, cids)
    return ID
\end{lstlisting}
\end{minipage}
    \caption{Pseudo code for adding an expression to the E-Graph. \texttt{find}
      looks for the given expression in the E-Graph, and returns its root ID if
      it already exists, or \texttt{NULL} otherwise. \texttt{insert} adds the
      given operator to the E-Graph, and points its children to E-Classes with the
      given class IDs. }
    \label{fig:egraphadd}
\end{figure}

The outer loop that matches equations to the graph can be implemented by a more
efficient algorithm like the Rete algorithm \cite{DBLP:journals/ai/Forgy82} when
the number of equations is large. However, we did not find matching to be
expensive and simply match by traversing the graph. Our implementation uses the
E-Graph data structure from the \texttt{egg}~\cite{EGG} library. 

\subsubsection*{Dealing with Expansive Rules}
While in theory equality saturation will converge with well-constructed rewrite
rules, in practice the E-Graph may explode for certain inputs under certain rules. For example, a long
chain of multiplication can be rewritten to an exponential number of
permutations under associativity and commutativity (AC rules). If we 
apply AC rules everywhere applicable in each iteration, the graph would soon use
up available memory. We call this application strategy the \emph{depth-first}
strategy because it eagerly applies expansive rules like AC. 
AC rules by themselves rarely affect performance, and
SystemML also provides the fused \texttt{mmchain} operator that efficiently computes
multiplication chains, so permuting a chain is likely futile. In practice, AC
rules are useful because they can enable other rewrites. Suppose we have a rule
$R_{factor}: A*X + B*X \rightarrow (A+B)*X$ and an expression $U*Y + Y*V$.
Applying commutativity to $Y*V$ would then transform the expression to be
eligible for $R_{factor}$. With this insight, we change each saturation
iteration to sample a limited number of matches to apply per rule, instead of
applying all matches. This amounts to adding \texttt{matches = sample(matches, limit)} between line 3 and line 4 in Figure~\ref{fig:egraphadd}. 
Sampling encourages each rule to be considered equally often
and prevents any single rule from exploding the graph. This helps ensure good
exploration of the search space when exhaustive search is impractical. 
But when it is
possible for saturation to converge and be exhaustive, it still converges with high probability
when we sample matches. Our experiments in Section~\ref{overhead} show sampling
always preserve convergence in practice.

\subsubsection*{Extracting the Optimal Plan}

\begin{figure}
\centering
    \includegraphics[scale=0.6]{CSE}
    \caption{The CSE problem. Each node shows its cost. A greedy optimizer picks nodes with costs 0, 1 and both nodes with cost 4; but the optimal plan uses nodes with costs 0, 2 and shares the same node with cost 4.}
    \label{cse}
\end{figure}

A greedy strategy to extract the best plan from the saturated E-Graph is to
traverse the graph bottom-up, picking the best plan at each level. This assumes
the best plan for any expression also contains the best plan for any of its
subexpressions. However, the presence of common subexpressions breaks this
assumption. In Figure~\ref{cse} each operator node is annotated with its cost. 
Between the nodes with costs 1 and 2, a greedy strategy would choose 1, which
incurs total cost of $1+4=5$. The greedy strategy then needs to pick the root
node with cost 0 and the other node with cost 4, incurring a total cost of 9. 
However, the optimal strategy is to pick the nodes with 0, 2 and share the same
node with cost 4, incurring a total cost of 6. 

We handle the complexity of the search problem with a constraint
solver. We assign a variable to each operator and each E-Class, then construct
constraints over the variables for the solver to select operators that make up a
valid expression. The solver will then optimize a cost function defined over the
variables; the solution then corresponds to the optimal expression equivalent to
the input.

\subsubsection*{Constraint Solving and Cost Function}
We encode the problem of extracting the cheapest plan from the E-Graph with
integer linear programming (ILP). Figure~\ref{ilp} shows this encoding. For each
operator in the graph, we generate a boolean variable $B_{op}$; for each E-Class
we generate a variable $B_c$. For the root class, we use the variable $B_r$.
Constraint $F(op)$ states that if the solver selects an operator, it must also
select all its children; constraint $G(c)$ states that if the solver selects an
E-Class, it must select at least one of its members. Finally, we assert $B_r$ must 
be selected, which constrains the extracted expression
to be in the same E-Class as the unoptimized expression. These three constraints together
ensure the
selected nodes form a valid expression equivalent to the unoptimized input.
Satisfying these constraints, the solver now minimizes the cost function given
by the total cost of the selected operators. Because each $B_{op}$ represents an
operator node in the E-Graph which can be shared by multiple parents, this
encoding only assigns the cost once for every shared common subexpression. In
our implementation, we use Gurobi~\cite{gurobi} to solve the ILP problem.

\begin{figure}
    \begin{align*}
Constraints &\equiv B_r \wedge \bigwedge_{op} F(op) \wedge \bigwedge_c G(c)
\\ F(op) &\equiv B_{op} \rightarrow \bigwedge_{c \in op.children} B_c \\ G(c)
&\equiv B_c \rightarrow \bigvee_{op \in c.nodes} B_{op}
\end{align*}
\[\textbf{minimize} \sum_{op} B_{op} \cdot C_{op} \textbf{ s.t. } Constraints\]

    \caption{ILP constraint and objective for extraction. }
    \label{ilp}
\end{figure}

Each operation usually has cost proportional to the output size in terms of
memory allocation and computation. Since the size of a matrix is proportional
its the number of non-zeroes (nnz), we use SystemML's estimate of nnz as the
cost for each operation. Under our relational interpretation, this corresponds
to the cardinality of relational queries. We use the simple estimation scheme in
Figure~\ref{fig:cost}, which we find to work well. Future work can hinge on the
vast literature on sparsity and cardinality estimation to improve the cost
model.

\begin{figure}
\begin{align*}
    \textbf{S}[X*Y] &= min(\textbf{S}[X], \textbf{S}[Y]) \\ \textbf{S}[X+Y] &=
    min(1, \textbf{S}[X]+\textbf{S}[Y]) \\ \textbf{S}[\sum_i X] &= min(1, |i|
    \cdot \textbf{S}[X])
\end{align*}
    \caption[Sparsity estimation]{Sparsity estimation. We define $sparsity = nnz
      / size$, i.e. a 0 matrix has sparsity $0.0$\protect\footnotemark. $|i|$ is
      the size of the aggregated dimension. }
    \label{fig:cost}
\end{figure}
\footnotetext{Some may find this definition counter-intuitive; we define it so
  to be consistent with SystemML.}

\subsection{Schema and Sparsity as Class Invariant}
In the rules $R_{EQ}$ used by the saturation process, Rule~(\ref{RRC_ma}) If $i
\not\in A$, $A * \sum_i B = \sum_i (A * B)$ contains a condition on attribute
$i$ which may be deeply nested in the expression. This means the optimizer
cannot find a match with a simple pattern match. Fortunately, all expressions in
the same class must contain the same set of free attributes (attributes not
bound by aggregates). In other words, the set of free variables is invariant
under equality. This corresponds precisely to the schema of a database -
equivalent queries must share the same schema. We therefore annotate each class
with its schema, and also enable each equation to match on the schema.

In general, we find class invariants to be a powerful construct for programming
with E-Graphs. For each class we track as class invariant if there is a constant
scalar in the class. As soon as all the children of an operator are found to
contain constants, we can fold the operator with the constant it computes. This
seamlessly integrates constant folding with the rest of the rewrites. We also
treat sparsity as a class invariant and track it throughout equality saturation.
Because our sparsity estimation is conservative, equal expressions that use
different operators may have different estimates. But as soon as we identify
them as equal, we can merge their sparsity estimates by picking the tighter one,
thereby improving our cost function. Finally, we also take advantage of the
schema invariant during constraint generation. Because we are only interested in
RA expressions that can be translated to LA, we only generate symbolic variables
for classes that have no more than two attributes in their schema. This prunes
away a large number of invalid candidates and helps the solver avoid wasting
time on them. We implement class invariants using \texttt{egg}'s Metadata API.

\subsection{Translation, Operator Fusion and Custom Functions}
Since equality saturation can rewrite any expression given a set of equations,
we can directly perform the translation between LA and RA within saturation,
simply by adding the translation rules $R_{LR}$ from Figure~\ref{RMR}.
Furthermore, saturation has flexible support for custom functions. The simplest
option is to treat a custom functions as a black box, so saturation can still
optimize below and above them. With a little more effort, we have the option to
extend our equations $R_{EQ}$ to reason about custom functions, removing the
optimization barrier. We take this option for common operators that are not part
of the core RA semantics, e.g. square, minus and divide. In the best scenario,
if the custom function can be modeled by a combination of basic operators, we
can add a rule equating the two, and retain both versions in the same graph for
consideration. In fact, this last option enables us to encode fused operators
and seamlessly integrate fusion with other rewrite rules. As a result, the
compiler no longer need to struggle with ordering fusion and rewrites, because
saturation simultaneously considers all possible ordering.

\subsection{Saturation v.s. Heuristics}

Using equality saturation, SPORES elegantly remedies the drawbacks of
heuristics mentioned in the beginning of section~\ref{explore}. First, when two
or more conflicting rewrites apply, they would be added to the same E-Class, and
the extraction step will pick the more effective one based on the global cost
estimate. Second, there is no need to carefully order rewrites, because
saturation simultaneously considers all possible orders. For example, when rules
$R_1$ and $R_2$ can rewrite expression $e$ to either $R_1(R_2(e))$ or
$R_2(R_1(e))$, one iteration of saturation would add $R_1(e)$ and $R_2(e)$ to
the graph, and another iteration would add both $R_1(R_2(e))$ and $R_2(R_1(e))$
to the same E-Class. Third, rules do not need to reason about their dependency on
input or program properties, because extraction uses a global cost model that
holistically incorporates factors like input sparsity and common subexpressions.
Finally, every rule application in saturation applies one step of rewrite on top
of those already applied, naturally composing complex rewrites out of simple
ones.



\subsection{Integration within SystemML}

We integrate SPORES into SystemML to leverage its compiler
infrastructure. Figure~\ref{warp} shows the architecture of the integrated
system: the optimizer plugs into the algebraic rewrite pass in SystemML. It
takes in a DAG of linear algebra operations, and outputs the optimized DAG.
Within the optimizer, it first translates the LA DAG into relational algebra,
performs equality saturation, and finally translates the optimal expression back
into LA. We obtain matrix characteristics such as dimensions and sparsity
estimation from SystemML. Since we did not focus our efforts in supporting
various operators and data types unrelated to linear algebra computation (e.g.
string manipulation), we only invoke SPORES on important LA expressions
from the inner loops of the input program.

\begin{figure}
    \centering \includegraphics[width=0.6\linewidth]{arch}
    \caption{Architecture of SPORES and integration within SystemML}
    \label{warp}
\end{figure}{}